% ============================================================
\chapter{Pricing des Options}
\label{ch:pricing}
% ============================================================

\section{Introduction}

La formule de Black-Scholes est \'el\'egante, mais elle ne s'applique qu'aux options europ\'eennes sur un sous-jacent suivant un mouvement brownien g\'eom\'etrique. Que faire pour une option am\'ericaine (exercice anticip\'e), une option asiatique (d\'ependant de la moyenne du cours), ou un produit exotique \`a barri\`ere~? Les formules ferm\'ees n'existent plus, et il faut recourir au calcul num\'erique. Trois grandes familles de m\'ethodes se partagent le terrain~: les arbres binomiaux, h\'erit\'es du mod\`ele de Cox-Ross-Rubinstein (1979), qui discr\'etisent le temps et construisent l'option par r\'ecurrence arri\`ere~; les m\'ethodes de Monte Carlo, qui simulent des milliers de trajectoires du sous-jacent et moyennent les payoffs actualis\'es~; et les diff\'erences finies, qui r\'esolvent num\'eriquement l'EDP de Black-Scholes sur une grille.

Chacune a ses forces et ses faiblesses, et le choix d\'epend du produit, de la dimension et de la pr\'ecision requise.

\section{Arbres binomiaux}

\subsection{Algorithme CRR}

Rappelons le mod\`ele de Cox--Ross--Rubinstein (chapitre 1). L'algorithme
proc\`ede en deux phases~:
\begin{enumerate}
    \item \textbf{Phase forward}~: construction de l'arbre des prix
          $S_{n,j} = S_0 u^j d^{n-j}$ pour $n = 0, \ldots, N$ et
          $j = 0, \ldots, n$.
    \item \textbf{Phase backward}~: calcul r\'ecursif des valeurs de
          l'option de $n = N$ \`a $n = 0$.
\end{enumerate}

\begin{formules}[title=\textbf{R\'ecurrence backward}]
Pour une option europ\'eenne~:
\[
    V_{n,j} = e^{-r\Delta t}\bigl[q \cdot V_{n+1,j+1}
              + (1-q) \cdot V_{n+1,j}\bigr],
\]
avec condition terminale $V_{N,j} = h(S_{N,j})$.

Pour une option am\'ericaine, on ajoute la possibilit\'e d'exercice
anticip\'e~:
\[
    V_{n,j} = \max\!\bigl(h(S_{n,j}),\;
              e^{-r\Delta t}[q V_{n+1,j+1} + (1-q) V_{n+1,j}]\bigr).
\]
\end{formules}

\subsection{Arbre trinomial}

\begin{definition}[Arbre trinomial]
L'arbre trinomial utilise trois mouvements possibles~:
\[
    S_{n+1} \in \{u S_n,\; S_n,\; d S_n\}
\]
avec des probabilit\'es $(p_u, p_m, p_d)$ choisies pour respecter
les deux premiers moments~:
\begin{align}
    p_u &= \frac{1}{2}\left(\frac{\sigma^2\Delta t + \nu^2\Delta t^2}
           {(\Delta x)^2} + \frac{\nu\Delta t}{\Delta x}\right), \\
    p_d &= \frac{1}{2}\left(\frac{\sigma^2\Delta t + \nu^2\Delta t^2}
           {(\Delta x)^2} - \frac{\nu\Delta t}{\Delta x}\right), \\
    p_m &= 1 - p_u - p_d,
\end{align}
o\`u $\nu = r - \sigma^2/2$ et $\Delta x = \sigma\sqrt{3\Delta t}$.
\end{definition}

\begin{remarque}
L'arbre trinomial converge plus rapidement que le binomial et pr\'esente
moins d'oscillations dans le prix en fonction de $N$.
\end{remarque}

\section{M\'ethodes de Monte Carlo}

\subsection{Principe g\'en\'eral}

\begin{theoreme}[Pricing par Monte Carlo]
Le prix d'un d\'eriv\'e europ\'een de payoff $h(S_T)$ est~:
\[
    V_0 = e^{-rT}\E^\mathbb{Q}[h(S_T)]
        \approx \frac{e^{-rT}}{M}\sum_{m=1}^{M} h(S_T^{(m)}),
\]
o\`u les $S_T^{(m)}$ sont des r\'ealisations ind\'ependantes de $S_T$
sous $\mathbb{Q}$.
\end{theoreme}

La convergence est en $O(1/\sqrt{M})$ par le th\'eor\`eme central limite,
ind\'ependamment de la dimension du probl\`eme. C'est l'avantage majeur
de Monte Carlo pour les probl\`emes en grande dimension.

\begin{proposition}[Intervalle de confiance]
L'estimateur Monte Carlo $\hat{V}_M$ satisfait~:
\[
    \PP\!\left(\abs{\hat{V}_M - V_0} \leq z_{\alpha/2}
    \frac{\hat{\sigma}_{\text{MC}}}{\sqrt{M}}\right) \approx 1 - \alpha,
\]
o\`u $\hat{\sigma}_{\text{MC}}$ est l'\'ecart-type empirique des payoffs
actualis\'es et $z_{\alpha/2}$ le quantile de la loi normale.
\end{proposition}

\subsection{Simulation du MBG}

Sous la mesure risque-neutre~:
\[
    S_T = S_0 \exp\!\left[\left(r - \frac{\sigma^2}{2}\right)T
          + \sigma\sqrt{T}\,Z\right], \quad Z \sim \mathcal{N}(0,1).
\]

Pour une simulation sur une grille $0 = t_0 < t_1 < \cdots < t_N = T$
(n\'ecessaire pour les options path-dependent)~:
\[
    S_{t_{i+1}} = S_{t_i} \exp\!\left[\left(r - \frac{\sigma^2}{2}\right)
    \Delta t + \sigma\sqrt{\Delta t}\,Z_i\right].
\]

\begin{codeblock}[title=\textbf{Monte Carlo pour un call europ\'een}]
\begin{minted}{python}
import numpy as np

def mc_european(S0, K, T, r, sigma, M=100000, option='call'):
    """Prix Monte Carlo d'une option europeenne."""
    Z = np.random.standard_normal(M)
    ST = S0 * np.exp((r - 0.5*sigma**2)*T + sigma*np.sqrt(T)*Z)

    if option == 'call':
        payoffs = np.maximum(ST - K, 0)
    else:
        payoffs = np.maximum(K - ST, 0)

    disc_payoffs = np.exp(-r*T) * payoffs
    price = np.mean(disc_payoffs)
    std_err = np.std(disc_payoffs) / np.sqrt(M)
    ci_95 = (price - 1.96*std_err, price + 1.96*std_err)
    return price, std_err, ci_95

np.random.seed(42)
price, se, ci = mc_european(100, 100, 1.0, 0.05, 0.20)
print(f"Prix MC: {price:.4f} +/- {se:.4f}")
print(f"IC 95%: [{ci[0]:.4f}, {ci[1]:.4f}]")
\end{minted}
\end{codeblock}

\subsection{R\'eduction de variance}

\begin{definition}[Variables antith\'etiques]
Pour chaque $Z_i$, on calcule aussi le payoff avec $-Z_i$, puis on
moyenne~:
\[
    \hat{V} = \frac{e^{-rT}}{M}\sum_{m=1}^{M}
    \frac{h(S_T^{(Z_m)}) + h(S_T^{(-Z_m)})}{2}.
\]
\end{definition}

\begin{definition}[Variable de contr\^ole]
On utilise un estimateur $Y$ de moyenne connue $\E[Y]$ corr\'el\'e au
payoff $X = h(S_T)$~:
\[
    \hat{V}_{\text{CV}} = \frac{1}{M}\sum_{m=1}^{M}
    \bigl[X^{(m)} - \beta(Y^{(m)} - \E[Y])\bigr],
\]
o\`u $\beta = \Cov(X,Y)/\Var(Y)$ est le coefficient optimal.
Un choix classique est $Y = S_T$ avec $\E^\mathbb{Q}[S_T] = S_0 e^{rT}$.
\end{definition}

\begin{codeblock}[title=\textbf{Monte Carlo avec r\'eduction de variance}]
\begin{minted}{python}
def mc_antithetic(S0, K, T, r, sigma, M=100000):
    """Monte Carlo avec variables antithetiques."""
    Z = np.random.standard_normal(M)
    ST_plus = S0 * np.exp((r - 0.5*sigma**2)*T + sigma*np.sqrt(T)*Z)
    ST_minus = S0 * np.exp((r - 0.5*sigma**2)*T - sigma*np.sqrt(T)*Z)
    payoffs = 0.5*(np.maximum(ST_plus - K, 0)
                   + np.maximum(ST_minus - K, 0))
    disc = np.exp(-r*T) * payoffs
    return np.mean(disc), np.std(disc)/np.sqrt(M)

def mc_control_variate(S0, K, T, r, sigma, M=100000):
    """Monte Carlo avec variable de controle (S_T)."""
    Z = np.random.standard_normal(M)
    ST = S0 * np.exp((r - 0.5*sigma**2)*T + sigma*np.sqrt(T)*Z)
    X = np.exp(-r*T) * np.maximum(ST - K, 0)
    Y = np.exp(-r*T) * ST
    EY = S0  # E^Q[exp(-rT)*S_T] = S_0

    beta = np.cov(X, Y)[0, 1] / np.var(Y)
    X_cv = X - beta * (Y - EY)
    return np.mean(X_cv), np.std(X_cv)/np.sqrt(M)

np.random.seed(42)
p1, s1 = mc_antithetic(100, 100, 1.0, 0.05, 0.20)
p2, s2 = mc_control_variate(100, 100, 1.0, 0.05, 0.20)
print(f"Antithetique: {p1:.4f} (se={s1:.4f})")
print(f"Controle:     {p2:.4f} (se={s2:.4f})")
\end{minted}
\end{codeblock}

\section{Diff\'erences finies}

L'\'EDP de Black--Scholes peut \^etre r\'esolue num\'eriquement par
la m\'ethode des diff\'erences finies.

\subsection{Transformation et discr\'etisation}

On pose $x = \ln(S/K)$, $\tau = T - t$. L'EDP se transforme en~:
\[
    \frac{\partial u}{\partial\tau} = \frac{\sigma^2}{2}
    \frac{\partial^2 u}{\partial x^2}
    + \left(r - \frac{\sigma^2}{2}\right)\frac{\partial u}{\partial x} - ru.
\]

\begin{definition}[Sch\'ema explicite]
On discr\'etise sur une grille $(x_i, \tau_j)$~:
\[
    \frac{u_i^{j+1} - u_i^j}{\Delta\tau}
    = \frac{\sigma^2}{2}\frac{u_{i+1}^j - 2u_i^j + u_{i-1}^j}{(\Delta x)^2}
    + \nu\frac{u_{i+1}^j - u_{i-1}^j}{2\Delta x} - ru_i^j.
\]
Ce sch\'ema est \textbf{explicite} (la valeur au pas suivant se calcule
directement) mais soumis \`a une condition de stabilit\'e CFL.
\end{definition}

\begin{definition}[Sch\'ema implicite (Crank--Nicolson)]
Le sch\'ema de Crank--Nicolson moyenne les pas explicite et implicite~:
\[
    \frac{u_i^{j+1} - u_i^j}{\Delta\tau}
    = \frac{1}{2}\bigl[\mathcal{L}u_i^j + \mathcal{L}u_i^{j+1}\bigr],
\]
o\`u $\mathcal{L}$ est l'op\'erateur diff\'erentiel spatial. Ce sch\'ema
est inconditionnellement stable et d'ordre 2 en temps et en espace.
\end{definition}

\begin{codeblock}[title=\textbf{Diff\'erences finies --- Crank--Nicolson}]
\begin{minted}{python}
import numpy as np
from scipy.linalg import solve_banded

def fd_crank_nicolson(S0, K, T, r, sigma, Nx=200, Nt=500):
    """Prix d'un call europeen par Crank-Nicolson."""
    x_max = 5 * sigma * np.sqrt(T)
    dx = 2 * x_max / Nx
    dt = T / Nt
    x = np.linspace(-x_max, x_max, Nx + 1)

    nu = r - 0.5 * sigma**2
    alpha = 0.5 * dt * sigma**2 / dx**2
    beta_val = 0.5 * dt * nu / (2 * dx)

    # Condition terminale (call)
    u = np.maximum(K * (np.exp(x) - 1), 0)

    # Coefficients tridiagonaux
    a = -0.5 * (alpha - beta_val)  # sous-diag
    b = 1 + alpha + 0.5 * r * dt  # diag
    c = -0.5 * (alpha + beta_val)  # sur-diag

    a_e = 0.5 * (alpha - beta_val)
    b_e = 1 - alpha - 0.5 * r * dt
    c_e = 0.5 * (alpha + beta_val)

    # Matrice bande pour solve_banded
    AB = np.zeros((3, Nx + 1))
    AB[0, 1:] = c       # sur-diag
    AB[1, :] = b         # diag
    AB[2, :-1] = a       # sous-diag

    for j in range(Nt):
        rhs = a_e * np.roll(u, 1) + b_e * u + c_e * np.roll(u, -1)
        rhs[0] = 0  # condition aux limites
        rhs[-1] = K * (np.exp(x[-1]) - np.exp(-r * (j+1)*dt))
        AB[1, 0] = 1; AB[0, 1] = 0
        AB[1, -1] = 1; AB[2, -2] = 0
        u = solve_banded((1, 1), AB, rhs)
        AB[1, 0] = b; AB[0, 1] = c
        AB[1, -1] = b; AB[2, -2] = a

    # Interpolation pour S = S0
    x0 = np.log(S0 / K)
    price = np.interp(x0, x, u)
    return price

price_fd = fd_crank_nicolson(100, 100, 1.0, 0.05, 0.20)
print(f"Prix Crank-Nicolson: {price_fd:.4f}")
\end{minted}
\end{codeblock}

\section{Comparaison des m\'ethodes}

\begin{center}
\begin{tabular}{lccc}
    \toprule
    \textbf{Crit\`ere} & \textbf{Arbres} & \textbf{Monte Carlo}
    & \textbf{Diff. finies} \\
    \midrule
    Dimension & faible & \'elev\'ee & faible \\
    Am\'ericaines & oui & difficile & oui \\
    Path-dependent & non & oui & difficile \\
    Convergence & $O(1/N)$ & $O(1/\sqrt{M})$ & $O(h^2)$ \\
    Grecques & diff. finies & bump-and-reprice & directes \\
    \bottomrule
\end{tabular}
\end{center}

\section{M\'ethode de Longstaff--Schwartz (options am\'ericaines)}

\begin{theoreme}[Algorithme LSM]
Pour une option am\'ericaine, l'algorithme de Longstaff--Schwartz (2001)
estime la valeur de continuation par r\'egression sur des polyn\^omes
des prix simul\'es. \`A chaque pas de temps $t_n$ en backward~:
\begin{enumerate}
    \item Calculer le payoff d'exercice imm\'ediat $h(S_{t_n}^{(m)})$.
    \item R\'egresser les flux futurs actualis\'es sur des fonctions de base
          $\psi_k(S_{t_n}^{(m)})$ (e.g., polyn\^omes de Laguerre).
    \item Exercer si $h(S_{t_n}^{(m)}) > \hat{C}(S_{t_n}^{(m)})$
          (continuation estim\'ee).
\end{enumerate}
\end{theoreme}

\begin{codeblock}[title=\textbf{Longstaff--Schwartz pour put am\'ericain}]
\begin{minted}{python}
def lsm_american_put(S0, K, T, r, sigma, M=50000, N=100):
    """Put americain par Longstaff-Schwartz."""
    dt = T / N
    discount = np.exp(-r * dt)

    # Simulation des trajectoires
    Z = np.random.standard_normal((N, M))
    S = np.zeros((N + 1, M))
    S[0] = S0
    for i in range(N):
        S[i+1] = S[i] * np.exp((r - 0.5*sigma**2)*dt
                                + sigma*np.sqrt(dt)*Z[i])

    # Payoff terminal
    cashflow = np.maximum(K - S[N], 0)
    stopping = np.full(M, N)

    # Backward induction
    for n in range(N - 1, 0, -1):
        exercise = np.maximum(K - S[n], 0)
        itm = exercise > 0

        if np.sum(itm) > 0:
            X = S[n, itm]
            Y = cashflow[itm] * discount**(stopping[itm] - n)
            # Regression polynomiale (degre 3)
            A = np.column_stack([np.ones_like(X), X, X**2, X**3])
            coeffs = np.linalg.lstsq(A, Y, rcond=None)[0]
            continuation = A @ coeffs

            exercise_now = exercise[itm] > continuation
            idx = np.where(itm)[0][exercise_now]
            cashflow[idx] = exercise[idx]
            stopping[idx] = n

    price = np.mean(cashflow * discount**stopping)
    return price

np.random.seed(42)
p_am = lsm_american_put(100, 100, 1.0, 0.05, 0.20)
print(f"Put americain (LSM): {p_am:.4f}")
\end{minted}
\end{codeblock}

\section{Exercices}

\begin{exercice}[Convergence des m\'ethodes]
\begin{enumerate}
    \item Comparer les prix CRR, Monte Carlo et Crank--Nicolson pour un
          call europ\'een avec les param\`etres $S_0=100$, $K=100$, $T=1$,
          $r=0.05$, $\sigma=0.20$.
    \item Tracer l'erreur par rapport \`a Black--Scholes en fonction du
          param\`etre de discr\'etisation ($N$ ou $M$).
    \item Mesurer le temps de calcul pour chaque m\'ethode.
\end{enumerate}
\end{exercice}

\begin{exercice}[R\'eduction de variance]
\begin{enumerate}
    \item Impl\'ementer les variables antith\'etiques et la variable de
          contr\^ole.
    \item Comparer les erreurs standard pour $M = 10^4, 10^5, 10^6$.
    \item Calculer le facteur de r\'eduction de variance pour chaque
          technique.
\end{enumerate}
\end{exercice}

\begin{exercice}[Options am\'ericaines]
\begin{enumerate}
    \item Comparer le prix d'un put am\'ericain par l'arbre CRR, par
          Crank--Nicolson avec exercice anticip\'e, et par LSM.
    \item D\'eterminer la prime d'exercice anticip\'e
          $P_{\text{am}} - P_{\text{eu}}$ en fonction de $r$ et $\sigma$.
\end{enumerate}
\end{exercice}

\begin{exercice}[Grecques par Monte Carlo]
\begin{enumerate}
    \item Calculer $\Delta$ et $\Gamma$ par \textit{bump-and-reprice}~:
          $\Delta \approx (V(S+h) - V(S-h))/(2h)$.
    \item Calculer $\Delta$ par la m\'ethode des \textit{likelihood ratios}~:
          $\Delta = e^{-rT}\E\bigl[h(S_T)\frac{Z}{\sigma S_0\sqrt{T}}\bigr]$.
    \item Comparer les \'ecarts-types des deux estimateurs.
\end{enumerate}
\end{exercice}
